\documentclass[11pt]{letter}
\usepackage[margin=2.2cm,top=1.5cm,bottom=1.5cm]{geometry}
\usepackage{enumitem}
\setlist{itemsep=1pt,topsep=3pt,leftmargin=*}
\signature{Dr. Nagesh Chougule\\ {}on behalf of all authors}
\address{Department of Mechanical Engineering\\ {}COEP Technological University\\ {}Pune, India\\ {}nkc.mech@coeptech.ac.in}
\begin{document}
\begin{letter}{The Editor-in-Chief\\ Reliability Engineering \& System Safety}
\opening{Dear Editor,}
\setlength{\parskip}{0.4em}

We submit the manuscript ``Evidence, not verdicts: what 590 public PX4 crash logs can and cannot tell a drone fleet operator'' as a research article.

Flight logs are increasingly used to decide why a small drone crashed, who was responsible, and which aircraft to
inspect next. We tested those uses on the largest public archive of real autopilot logs (PX4 Flight Review, 479,547
logs). Each confirmatory test was committed to version control with its pass mark before it was run;
analyses added during review are marked as exploratory.

Two findings are, to our knowledge, new and apply to any evaluation of log-based diagnostics:
\begin{itemize}
\item \textbf{Evaluation choices create much of the apparent signal.} A factorial audit of 32 evaluation settings
attributes 0.21 [0.06, 0.34] of a crash-vs-normal AUC of 0.88 to choices of samples and
validation. Only one choice mattered on its own: grouping
cross-validation folds by vehicle, where drones contribute several flights. Comparing each drone with its own earlier flights brought the AUC to
chance.
\item \textbf{Flight order cannot be taken from archive dates.} GPS time showed that 40--50\% of crash-day
``earlier'' flights (47 of 95 in one sample, 85 of 210 in another) were flown after the crash. Our main results
held without them; an apparent within-drone effect did not, and we have withdrawn it.
\end{itemize}
These rest on declared tests:
\begin{itemize}
\item neither a declared rule nor cross-validated models separated pilot-caused from aircraft-caused crashes
(balanced accuracy 0.47 on 112 held-out crashes);
\item one rule-based check owed its accuracy to the crash itself; for models, leakage was immaterial;
\item the flight that crashed was no different from the same drone's earlier flights, whereas drones that later
crashed vibrated more than matched drones on ordinary flights, a suggestive signal ($p=0.012$ clustered by airframe,
Holm-adjusted 0.08) useful for inspection ranking, not crash warning;
\item small language-model panels did worse than a fixed rule and could be steered by one compromised
reviewer.
\end{itemize}
The findings feed a recommended reporting checklist and an open-source fleet-assurance tool, AeroFleet~2.0, that
presents evidence rather than verdicts. Code, declarations, derived features and labels are released.

The manuscript is original, has not been published and is not under consideration elsewhere. All authors approve the
submission. An AI assistant was used in preparing code and text, as declared in the manuscript.

\closing{Yours sincerely,}
\end{letter}
\end{document}
